\scanpage{59}
cubique \inserted{$(\unsure{\mathscr T|\Gamma|\mathcal T}_2)$} muni d'un $\mathscr T_1$ et d'un $\mathscr T_0$ contenu dans celui-ci, équivaut à la \deleted{donnée \unsure{d'un}} \inserted{catégorie} \deleted{de} $(\mathrm{Ens})_{3,3}\times(\mathrm{Ens})_2$ \deleted{ens à trois éléments,} \deleted{en associant à un} \inserted{[en associant l'ens des} triades de $X$, et l'ens des deux orbites dans $S-\unsure{\mathscr X|\mathfrak X|\mathcal X}$]

\deleted{couple de} La \deleted{donnée} \inserted{catégorie} des systèmes $(S\supset\mathscr X\supset X\supset t)$ d'un $\mathscr T_2$ contenant un $\mathscr T_1$ contenant un $\mathscr T_0$ contenant un $\mathscr T_{-1}$, équivaut à la catégorie $(\mathrm{Ens})_3\times(\mathrm{Ens})_2\times(\mathrm{Ens})_2$, [en lui associant l'ens des points de l'une des triades de $\deleted{\unsure{X|\mathscr X|\mathcal X}}$ pour la relation d'équiv. définie par $t$, l'ens des deux coorientations de $X$, et l'ens des deux orbites de $S-\mathscr X$] ....

\textcircled{17} Soit $\mathscr X$ un $\mathscr T_1$ dans \deleted{\unsure{$S$}}\unsure{un|le} $\mathscr T_2$ $S$. Alors $S-\mathscr X$ est une bibase de $S$, et on trouve ainsi une bijection entre l'ens des $\mathscr T_1$ de $S$, et l'ens des bibases.
La catégorie des $\mathscr T_2$ munis d'une bibase, celle des $\mathscr T_2$ munis d'un $\unsure{\mathscr X_1|\mathcal X_1}$, et la catégorie des $(\mathrm{Ens})_6\times(\mathrm{Ens})_2$, sont canoniquement équivalentes.

% archive sheet number: 58

